% !TeX spellcheck = ru_RU
% !TEX root = vkr.tex

\section{Структура учебного плана}
\label{sec:relatedworks}

Учебный план является одним из основным документов, определяющих содержание образовательной программы, в частности \enquote{перечень, трудоемкость, последовательность и распределение по периодам обучения учебных предметов, курсов, дисциплин (модулей), практики, иных видов учебной деятельности и [...] формы промежуточной аттестации обучающихся}~\cite[п. 22]{Statya2Osnovnye}.

В СПбГУ его содержание описывается не только федеральным государственным образовательным стандартом (например,~\cite{PrikazMinisterstvaObrazovaniya}), но и собственным образовательным стандартом СПбГУ~\cite{ObrazovatelnyeStandartyTrebovaniya}.

Таким образом, в соответствии с требованиями нормативных актов, модель учебного плана в информационных системах СПбГУ выглядит примерно так.

Высшей сущностью является сам учебный план, содержащий информацию вида языка реализации, формы обучения и подобного, а также семестры и их содержимое.
Каждый семестр состоит из трех частей:
\begin{enumerate}
    \item Базовая часть.
    \item Вариативная часть.
    \item Факультативная часть.
\end{enumerate}

Стоит отметить, что вариативная часть в текущей редакции образовательного стандарта не является обязательной\footnote{Заметка о разнице ФГОС ВО 3+ и ФГОС ВО 3++: \url{https://vk.com/@metodist1000-kak-opredelit-fgos-vo-3-ili-3} (дата обращения: 11.02.2026).}, и расположение предметов в базовой
или вариативной части~--- методический, а не формальный вопрос.
Например, дисциплина \enquote{Английский язык} в действующих учебных планах, как правило, относится к базовой части, хотя согласно предыдущим версиям стандартов была бы отнесена к вариативной части.

Следующей единицей дробления является блок предметов, их можно разделить на два вида, так называемые \enquote{простые} блоки и \enquote{сложные} блоки.

\enquote{Простой} блок, также часто называемые \enquote{дисциплина по выбору}, типичный способ выразить элективные дисциплины.
Также через \enquote{простые} блоки выражаются одиночные дисциплины, в таком случае блок просто содержит ровно одну дисциплину.

\enquote{Сложный} блок, также называемый \enquote{блок дисциплин по выбору}, позволяет выстроить \enquote{трек} сквозь несколько семестров, в таком случае требуется и название набору \enquote{треков}, и каждому конкретному \enquote{треку}.
При этом внутри \enquote{сложного} блока естественным образом содержатся \enquote{простые} блоки.
